\subsection*{一步一步建立直观}

\paragraph{1. 把多项式做成局部仿制品}

首先让近似函数在中心点拥有相同高度，再让斜率相同，然后匹配弯曲及更高阶变化。系数除以阶乘，是因为对幂函数重复求导会产生相应的阶乘。右图比较指数函数与在零点展开的多项式，每加一阶就增加一个局部匹配条件。



\paragraph{2. 近似好坏同时取决于阶数和距离}

固定阶数而把观察点远离中心，误差通常变大。右图的数值读数比肉眼重合更可靠，但仍只是一个点的误差；要承诺整段区间的精度，还须对余项给出统一上界。近似式如果不附中心与误差范围，很容易被使用到不适合的地方。



\paragraph{3. 区分有限公式与无穷级数}

增加阶数，在同一点比较误差是否缩小。指数函数的泰勒级数在每个实数点收敛到原函数，但这并不代表任意无穷可导函数都等于自己的泰勒级数。有限阶公式配合余项先成立；让阶数趋向无穷，还需要余项确实趋零。



\subsection*{定义、命题与证明}

\paragraph{定义：指数函数在零点的泰勒多项式}

定义 $P_n(x)=\sum_{k=0}^n x^k/k!$。每一项用于匹配指数函数在零点的一个更高阶导数。

\paragraph{命题：在固定区间内控制截断误差}

若 $|x|\le R$，则 $|e^x-P_n(x)|\le e^R R^{n+1}/(n+1)!$。此界对区间中的所有点同时成立。

\begin{proof}

由带拉格朗日余项的泰勒公式，余项为 $e^\xi x^{n+1}/(n+1)!$，其中 $\xi$ 位于零和 $x$ 之间。指数函数具有公式所需的各阶导数。

因为 $|\xi|\le R$，有 $e^\xi\le e^R$；再用 $|x|^{n+1}\le R^{n+1}$ 得到上界。图中单点实际误差通常更小；统一上界的作用是事先保证整个使用区间的精度。

\end{proof}
